% ==============================================================
%  GLM-4.6 — Model Card
%  Compiles on Overleaf with pdfLaTeX. One file, nothing to upload.
%  Export as:  GLM-4.6-Model-Card.pdf
% ==============================================================
\documentclass[11pt,a4paper]{article}

\usepackage[T1]{fontenc}
\usepackage[utf8]{inputenc}
\usepackage{XCharter}                       % warm serif for body text
\usepackage[scaled=0.92]{sourcesanspro}     % clean sans for headings
\usepackage{inconsolata}                    % mono

\usepackage[margin=2.6cm,top=2.8cm,bottom=2.8cm]{geometry}
\usepackage{xcolor}
\usepackage{titlesec}
\usepackage{enumitem}
\usepackage{booktabs}
\usepackage{tabularx}
\usepackage{array}
\usepackage[most]{tcolorbox}
\usepackage{fancyhdr}
\usepackage[hidelinks]{hyperref}

% ---------------- palette ----------------
\definecolor{paper}{HTML}{FAF9F5}
\definecolor{ink}{HTML}{191917}
\definecolor{muted}{HTML}{6B6760}
\definecolor{rule}{HTML}{DDD8CE}
\definecolor{clay}{HTML}{BF5B3D}
\definecolor{moss}{HTML}{4C6B3C}
\definecolor{stone}{HTML}{8A857C}

\pagecolor{paper}
\color{ink}

\hypersetup{
  pdftitle={GLM-4.6 — Model Card},
  pdfsubject={Generative AI coursework},
  pdfkeywords={GLM-4.6, Zhipu AI, model card, mixture of experts, open weights}
}

% ---------------- evidence tags ----------------
\newcommand{\tg}[2]{{\scriptsize\sffamily\bfseries\color{#1}[\,#2\,]}}
\newcommand{\Conf}{\tg{moss}{CONFIRMED}}
\newcommand{\Infr}{\tg{clay}{INFERRED}}
\newcommand{\Undi}{\tg{stone}{NOT DISCLOSED}}

% ---------------- headings ----------------
\titleformat{\section}
  {\sffamily\large\bfseries\color{ink}}
  {\color{clay}\thesection}{0.7em}{}
\titleformat{\subsection}
  {\sffamily\normalsize\bfseries\color{ink}}{}{0em}{}
\titlespacing*{\section}{0pt}{22pt}{7pt}
\titlespacing*{\subsection}{0pt}{14pt}{4pt}

\setlist[itemize]{leftmargin=1.3em,itemsep=5pt,topsep=5pt,parsep=0pt,
                  label=\textcolor{rule}{\rule[0.24em]{0.45em}{0.7pt}}}
\renewcommand{\arraystretch}{1.28}
\setlength{\parindent}{0pt}
\setlength{\parskip}{7pt}

% ---------------- a quiet callout: left rule, no fill ----------------
\newtcolorbox{aside}{
  enhanced, breakable, sharp corners, boxrule=0pt,
  colback=paper, colframe=paper,
  borderline west={1.6pt}{0pt}{clay},
  left=14pt, right=6pt, top=5pt, bottom=5pt,
  fontupper=\small\color{muted}
}

% ---------------- page furniture ----------------
\pagestyle{fancy}\fancyhf{}
\renewcommand{\headrulewidth}{0pt}
\fancyfoot[C]{\sffamily\footnotesize\color{stone}\thepage}

\newcommand{\thd}[1]{{\sffamily\bfseries\footnotesize\color{muted}#1}}

\begin{document}

% ==============================================================
%  TITLE
% ==============================================================
\thispagestyle{empty}
\vspace*{-0.4cm}

{\sffamily\footnotesize\bfseries\color{clay} MODEL CARD}

\vspace{12pt}
{\sffamily\fontsize{30}{34}\selectfont\bfseries\color{ink} GLM-4.6}

\vspace{9pt}
{\fontsize{13}{18}\selectfont\color{muted}
An open-weight mixture-of-experts language model from Zhipu AI, documented from public
sources --- with every claim marked by how well it is actually sourced.}

\vspace{16pt}
{\small\color{muted}
\textbf{\color{ink}Disha Kataria} (16014223034) \quad and \quad \textbf{\color{ink}Jonan Johnson} (16014223042)\\[3pt]
Generative AI coursework \quad\textbullet\quad August 2026}

\vspace{15pt}
{\color{rule}\hrule height 0.8pt}

% ==============================================================
\section*{At a glance}

\begin{tabularx}{\textwidth}{@{}>{\raggedright\arraybackslash}p{3.9cm}X@{}}
\toprule
\thd{Developer}      & Zhipu AI, operating publicly as Z.ai \\
\thd{Released}       & Late September 2025 \\
\thd{Architecture}   & Mixture-of-Experts Transformer, 355\,B total parameters \\
\thd{Context window} & 200K tokens in, up to 128K tokens out \\
\thd{License}        & MIT --- weights downloadable and commercially usable \\
\thd{Built for}      & Agentic task execution, software coding, long-context reasoning \\
\thd{Predecessor}    & GLM-4.5 \\
\bottomrule
\end{tabularx}

\subsection*{How to read this card}

Zhipu AI published a good deal about GLM-4.6 and stayed quiet about a good deal more. Rather
than smooth over the difference, we tagged every substantive claim:

\begin{itemize}[itemsep=4pt]
  \item \Conf\ \ Stated directly in the developer's own release material.
  \item \Infr\ \ Reasoned from the documented GLM family pattern. Not stated outright for
        GLM-4.6, and not safe to cite as a specification.
  \item \Undi\ \ Never published. We left the field empty rather than borrowing a number from a
        comparable model.
\end{itemize}

The tags are the honest part of this card. A fair share of what a complete model card should
contain is simply not in the public record, and we thought that was worth showing rather than
hiding.

% ==============================================================
\section{Purpose}

GLM-4.6 is a general-purpose large language model, part of Zhipu AI's General Language Model
series. It generates and reasons over natural language text for a broad range of downstream
applications, with particular emphasis on three areas: \textbf{agentic task execution} (tool use
and multi-step planning), \textbf{software coding} (generation, editing, and terminal-based
tasks), and \textbf{long-context reasoning}. \Conf

It is positioned as an open-weight alternative to proprietary frontier models. Because the
weights are downloadable under an MIT license, enterprises and researchers can self-host,
fine-tune, and deploy it without depending on a third-party API. The stated goal is to narrow
the gap between open and closed models on real-world coding and agentic benchmarks while staying
fully inspectable. \Conf

% ==============================================================
\section{Dataset}

Zhipu AI has not published a detailed breakdown of GLM-4.6's pretraining corpus. This is normal
practice for frontier-scale models, and we treat it as a transparency limitation rather than an
oversight --- see Section~6. From the developer's technical reports and the GLM family's
documented lineage, three things can be said with reasonable confidence:

\begin{itemize}
  \item The pretraining corpus is large-scale and multilingual, primarily Chinese and English
        web text, supplemented with code repositories, technical and academic documents, and
        dialogue data --- consistent with GLM-4 and GLM-4.5. \Infr
  \item GLM-4.6 builds on the GLM-4.5 architecture and training pipeline, with additional
        training extending context handling to 200K tokens and refreshing coding- and
        agent-relevant data: tool-use trajectories, terminal sessions, and front-end UI
        generation examples. \Infr
  \item Post-training data includes curated instruction-following, coding, and agentic-workflow
        datasets for supervised fine-tuning, plus preference and reward data for reinforcement
        learning. \Conf
\end{itemize}

Four fields that a complete dataset section would carry are absent from the public record
entirely: \textbf{total corpus size}, \textbf{source distribution and weighting},
\textbf{filtering and deduplication method}, and \textbf{data licensing status}. \Undi

\begin{aside}
We looked for these and could not find them. If you come across a specific figure for GLM-4.6's
dataset size in a blog post or comparison table, it is an estimate someone made --- not a number
Zhipu AI released.
\end{aside}

% ==============================================================
\section{Training method}

The base architecture is a Mixture-of-Experts Transformer with 355\,B total parameters, of which
only a subset of experts activates per token --- which is how a model this large stays
affordable to run. \Conf

The training pipeline, as far as it can be reconstructed, runs in four stages:

\begin{itemize}
  \item \textbf{Pretraining.} Multi-stage training on the corpus described in Section~2. \Infr
  \item \textbf{Context extension.} A dedicated phase stretching the input window to 200K
        tokens. \Infr
  \item \textbf{Supervised fine-tuning.} Curated instruction, coding, and agentic-task data,
        teaching the base model to follow instructions and use tools reliably. \Conf
  \item \textbf{Reinforcement learning.} Tuned for coding accuracy, agentic task completion, and
        reasoning stability. Zhipu AI reports this improved token efficiency by roughly
        \textbf{15\%} --- about 15\% fewer tokens to finish comparable tasks than GLM-4.5
        took. \Conf
\end{itemize}

Compute is a blank. Hardware, total training FLOPs, and training duration were never published,
so we have omitted them rather than estimating from parameter count. \Undi

% ==============================================================
\section{Evaluation}

Zhipu AI reports GLM-4.6's performance through its own benchmark suites rather than against a
single certification standard. The main axes are:

\begin{itemize}
  \item \textbf{CC-Bench}, an internally developed, human-evaluated real-world coding benchmark.
        GLM-4.6 is reported to reach near-parity with Claude Sonnet 4, and the developer's own
        materials acknowledge it still trails Claude Sonnet 4.5 on coding. \Conf
  \item \textbf{Token efficiency}, measured as tokens consumed per completed task --- roughly
        15\% fewer than GLM-4.5 for comparable work. \Conf
  \item \textbf{General capability} benchmarks covering reasoning, mathematics, and multilingual
        understanding, consistent with the broader GLM-4.x suite. Public GLM-4.6-specific numbers
        are thinner here than for the GLM-4.6V multimodal variant. \Infr
\end{itemize}

\medskip
\begin{tabularx}{\textwidth}{@{}>{\raggedright\arraybackslash}p{5.0cm}X>{\raggedleft\arraybackslash}p{2.2cm}@{}}
\toprule
\thd{Model} & \thd{Availability} & \thd{Win rate} \\
\midrule
Claude Sonnet 4.5 & Proprietary, Anthropic & 97\% \\
Claude Sonnet 4   & Proprietary, Anthropic & 90\% \\
\textbf{GLM-4.6}  & \textbf{Open weight, Zhipu AI} & \textbf{88\%} \\
GLM-4.5           & Open weight, Zhipu AI  & 76\% \\
\bottomrule
\end{tabularx}

\begin{aside}
These are the developer's own numbers, scored by evaluators the developer engaged, and no
independent lab has reproduced them. That does not make them wrong --- a company conceding that
a competitor beats its own model is not the behaviour of someone gaming a benchmark --- but a
comparative claim nobody outside the company has checked is still an unchecked claim.
\end{aside}

% ==============================================================
\section{Bias analysis}

There is no quantitative bias or fairness audit for GLM-4.6. No demographic parity testing, no
BBQ or StereoSet result, nothing equivalent --- from the developer or from anyone else. \Undi

That absence is the finding. What follows is inference from how the model was built, and we want
to be clear that these are \emph{reasoned risks, not measurements}:

\begin{itemize}
  \item \textbf{Language and cultural reach.} Training data drawn predominantly from Chinese- and
        English-language web sources will likely produce stronger, more culturally attuned output
        in those languages, and weaker, more error-prone output for lower-resource languages and
        the cultural contexts attached to them. \Infr
  \item \textbf{Inherited web-scrape skew.} As with essentially every large web-trained model,
        GLM-4.6 likely carries the representational imbalances of its sources --- some demographic
        groups, viewpoints, and regions over-represented, others thin on the ground --- which can
        surface as skewed associations in generated text. \Infr
  \item \textbf{What the tuning prioritised.} Post-training emphasised coding and tool use.
        Conversational and creative-writing bias behaviour has probably had less dedicated
        attention than it would in a model built primarily for chat. \Infr
\end{itemize}

\begin{aside}
If you are putting this model in front of users, run your own bias evaluation on your own
population first. That is ordinary good practice with any model; here it is the only evidence
that will exist.
\end{aside}

% ==============================================================
\section{Ethical considerations}

\begin{itemize}
  \item \textbf{The transparency gap.} Training data composition, red-teaming methodology, and
        the reinforcement learning reward criteria are not disclosed. This is what makes
        independent verification of any safety or fairness claim impossible, and it is the thread
        running through Sections~2, 4, and 5. \Undi
  \item \textbf{Open weights cut both ways.} A permissive MIT license with downloadable weights
        means whatever safety tuning ships with the model can be fine-tuned back out by anyone
        downstream. Openness is genuinely valuable, and it relocates responsibility for safe
        deployment onto whoever is hosting. \Conf
  \item \textbf{Agentic capability raises the stakes.} A model designed for tool use, terminal
        access, and multi-step execution can act on a mistake rather than merely state one.
        Production deployments warrant a human in the loop. \Infr
  \item \textbf{Jurisdiction is not inherited.} As a model from a China-based developer,
        deployments in some regions may need separate evaluation of data governance,
        export-control, and content-moderation requirements. \Infr
  \item \textbf{Say when it is AI.} Output used in academic, legal, or journalistic work should be
        disclosed as AI-generated or AI-assisted --- including, for the avoidance of doubt, work
        like this. \Conf
\end{itemize}

% ==============================================================
\section{Limitations}

\begin{itemize}
  \item \textbf{The documentation itself.} Dataset composition, training compute, and safety
        evaluation methodology are not fully disclosed, which caps how complete Sections~2, 4,
        and 5 of this card can be. \Undi
  \item \textbf{It is not the best coding model available.} Zhipu AI's own disclosures place
        GLM-4.6 behind the strongest proprietary coding models, Claude Sonnet 4.5 among
        them. \Conf
  \item \textbf{Open does not mean cheap.} At 355\,B parameters, even with sparse activation,
        self-hosting needs serious GPU infrastructure. The license is free; the hardware is
        not. \Conf
  \item \textbf{The usual LLM caveats.} Hallucination, sensitivity to how a prompt is phrased,
        uneven performance well outside the coding and agentic domain it was tuned for, and a
        fixed knowledge cutoff. \Infr
  \item \textbf{Nobody independent has checked.} At the time of writing we found no major
        third-party audit or standardised red-teaming report for GLM-4.6 in public sources. \Undi
\end{itemize}

% ==============================================================
\section{Appropriate use}

\subsection*{Well suited for}
\begin{itemize}
  \item Software development assistance --- code generation, debugging, refactoring, and terminal
        or agentic coding workflows.
  \item Front-end and UI generation, where the model has been specifically noted for polished
        output.
  \item Long-document analysis and summarisation, using the 200K-token context window.
  \item Self-hosted deployments needing data privacy, on-premise control, or fine-tuning on a
        proprietary codebase --- the case the MIT license exists to serve.
  \item Research into open-weight MoE behaviour and agentic AI systems.
\end{itemize}

\subsection*{Needs additional safeguards}
\begin{itemize}
  \item High-stakes autonomous decisions in medical, legal, or financial settings. No
        domain-specific safety evaluation has been published, so human review is not optional.
  \item Deployments serving low-resource languages or underrepresented cultural contexts, without
        evaluation of your own.
  \item Fully autonomous action with real-world consequences --- payments, system administration,
        irreversible file operations --- without a human in the loop.
  \item Anything requiring an audited fairness guarantee. This one cannot be satisfied at all
        right now, because no such audit is public.
\end{itemize}

% ==============================================================
\section*{Closing note}

GLM-4.6 is a genuinely open model, released under a license that means what it says, and
benchmarked candidly enough by its developer to admit a competitor is ahead. Those are real
virtues and they are not universal.

The gap is documentation. Data provenance, training compute, and fairness evaluation are the
three areas where an open-weight release most needs a paper trail --- because the license hands
capability to people who will run the model in contexts its developer never sees --- and those
are precisely the three areas left blank. The weights are open; the record behind them is not.

\vspace{10pt}
{\color{rule}\hrule height 0.8pt}
\vspace{6pt}

{\footnotesize\color{muted}
Prepared as an academic exercise from publicly available secondary sources --- developer blog
posts, technical press coverage, and published benchmark reports --- as of August 2026. Not an
official Zhipu AI publication. Where we could not find a primary source, we said so instead of
guessing.}

\end{document}
